\section{Методика тестирования и безопасность}

\subsection{Протокол безопасности полетов}

\subsubsection{1. Подготовка зоны (Pre-flight)}
\begin{itemize}
    \item Очистить зону $3 \times 3$ метра.
    \item Закрыть окна (ветер).
    \item Убрать животных и детей.
    \item Предупредить окружающих о запуске.
    \item Надеть защитные очки (для оператора).
\end{itemize}

\subsubsection{2. Проверка дрона (Checklist)}
\begin{itemize}
    \item \textbf{Пропеллеры:} Целые, затянуты, вращаются свободно.
    \item \textbf{Аккумулятор:} Заряжен > 3.8В на банку, закреплен надежно.
    \item \textbf{Камера:} Линза чистая, шлейф не болтается.
    \item \textbf{Wi-Fi:} Сигнал стабильный (RSS > -60 dBm).
    \item \textbf{Пионер Мини~2:} пульт включён, тумблер \texttt{SWB} в нижнем положении (автопилот проверяет наличие пилота), программа содержит блок \texttt{finally} с посадкой, прерывание \key{Ctrl}+\key{C} проверено.
\end{itemize}

\subsubsection{3. Экстренная ситуация}
Если дрон ведет себя неадекватно:
\begin{enumerate}
    \item Громко крикнуть <<ПОСАДКА!>>.
    \item Остановить программу и посадить дрон (кнопка LAND в наземной станции, если есть; у Мини~2 --- \key{Ctrl}+\key{C} в терминале, он ведёт в \texttt{finally} и к посадке, или кнопка \ui{Стоп} в Bricks); если он не садится --- выключить моторы тумблером взвода на пульте или командой Disarm.
    \item Если нет реакции — аккуратно поймать снизу (за батарею).
    \item Отключить аккумулятор.
\end{enumerate}

\subsection{Тестирование в разных условиях}

\subsubsection{Погодные условия (Ветер)}
Пионер --- лёгкий учебный дрон, и ветер сносит его заметно. Ориентир: не летать на улице при ветре сильнее 3 м/с.
Для теста в помещении используйте вентилятор:
\begin{itemize}
    \item Включите вентилятор на минимальную скорость с расстояния 2 метра.
    \item Попробуйте удержать точку (режим Optical Flow).
    \item Если дрон сносит > 50 см — условия непригодны для точной навигации.
\end{itemize}

\subsubsection{Освещение (Lux)}
Камера Pioneer имеет низкий динамический диапазон.
\begin{itemize}
    \item \textbf{Яркое солнце:} Пересвет (белое пятно). Маркеры не видны.
    \item \textbf{Сумерки:} Сильный шум матрицы. Маркеры <<дрожат>>.
    \item \textbf{Мерцание ламп:} Люминесцентные лампы (50 Гц) могут давать полосы на видео.
\end{itemize}

\subsection{Проверка в симуляторе перед полётом}
Прежде чем запускать программу на настоящем дроне, прогоните её в симуляторе --- так ловятся ошибки логики, а не железа.
\begin{itemize}
    \item \textbf{Несколько полигонов.} Запустите программу на двух-трёх seed процедурного полигона (миссии Python №\,55 и №\,62) и после сброса сцены. Алгоритм не должен опираться на координаты, которые совпали случайно (миссия №\,107).
    \item \textbf{Неуспешные ветки.} Проверьте исходы «метки нет» (несуществующий ID, как 249 в миссии №\,108), «кадр не получен» (\texttt{None}) и «точка не достигнута за отведённое время». Каждый должен заканчиваться понятным сообщением и посадкой.
    \item \textbf{Ограничения.} Скорости ограничены, высота --- с запасом над препятствиями, зона полёта задана границами.
    \item \textbf{Перенос.} Уберите расширения симулятора (\texttt{Camera(view="bottom")}), заново проверьте знаки поправок коротким импульсом, уменьшите скорости для первого полёта.
    \item \textbf{Мини~2.} Выберите модель \ui{Мини 2} и проверьте программу на \texttt{pioneer\_sdk2}. Симулятор упрощает или не поддерживает параметры автопилота, GPS, запись видео и фото, а нейросеть заменяет имитацией (главы~\ref{sec:sdk2} и~\ref{sec:rknn}): такие места проверяйте только на дроне и в первую очередь без пропеллеров.
\end{itemize}

\section{Готовые конфигурации (Templates)}

\subsection{Файл конфигурации camera\_params.yaml}
Те же параметры камеры, что сохраняют скрипты калибровки в \texttt{data.yml}, но в формате для сторонних инструментов: матрица --- \texttt{camera\_matrix}, дисторсия --- \texttt{distortion\_coefficients}. Примеры этой книги читают \texttt{data.yml}.
\begin{codefiletitle}[yaml]{camera_params.yaml}{python/advanced_navigation/camera_params_yaml.yaml}
\end{codefiletitle}

\subsection{Шаблон скрипта полета (main.py)}
Шаблон гарантирует посадку при любой ошибке: \texttt{land()} и пауза до касания выполняются в блоке \texttt{finally}, и только после них --- \texttt{disarm()}. Выключить моторы сразу после \texttt{land()} значило бы уронить дрон с высоты.
\begin{codefiletitle}[python]{Базовый шаблон}{python/advanced_navigation/main_template_py.py}
\end{codefiletitle}
